# Руководство пользователя PERFORM

PERFORM - платформа для управления тренировочным процессом. Она помогает тренеру вести спортсменов, планировать подготовку, назначать тренировки, контролировать готовность, фиксировать результаты и видеть аналитику. Спортсмен использует платформу для ежедневной готовности, просмотра назначенной работы и отметки фактического выполнения.

## 1. Роли пользователей

В платформе есть три основные роли.

| Роль | Что делает |
| --- | --- |
| Тренер | Ведёт спортсменов, создаёт планы, назначает тренировки, смотрит готовность, результаты и аналитику. |
| Спортсмен | Заполняет ежедневную готовность, смотрит назначенные тренировки, отмечает выполнение, фиксирует результаты. |
| Администратор | Управляет доступом и техническими настройками платформы. |

Основная работа идёт через верхнее меню. Если тренер выбрал спортсмена в верхнем меню, все рабочие вкладки показывают данные именно этого спортсмена.

## 2. Вход в систему

1. Откройте сайт платформы.
2. Введите email и пароль.
3. Нажмите **Войти**.
4. После входа откроется рабочая зона, доступная вашей роли.

Если пароль забыт, обратитесь к администратору или тренеру с правами управления пользователями.

## 3. Рабочие зоны тренера

У тренера доступны основные разделы:

| Раздел | Назначение |
| --- | --- |
| Dashboard | Быстрый обзор состояния спортсменов и важных сигналов. |
| Спортсмены | Список спортсменов, карточка профиля и данные спортсмена. |
| Планирование | План подготовки, мезоциклы, недельный план, шаблоны, сезон, календарь. |
| Аналитика | Тренды готовности, нагрузки, выполнения и рекомендации. |
| Review | Сравнение плана и факта выполнения. |

## 4. Выбор спортсмена тренером

1. В верхнем меню найдите поле **Спортсмен**.
2. Выберите нужного спортсмена.
3. После выбора все вкладки тренера работают с этим спортсменом.

Не нужно повторно выбирать спортсмена внутри каждой вкладки. Если где-то требуется назначение плана, система берёт спортсмена из верхнего меню.

## 5. Dashboard тренера

Dashboard нужен для быстрого контроля.

Что смотреть:

- текущую готовность спортсмена;
- активные назначения;
- ближайший старт;
- предупреждения;
- показатели выполнения;
- сигналы, требующие внимания тренера.

Используйте Dashboard как первый экран дня: сначала проверьте состояние спортсменов, затем переходите к планированию или review.

## 6. Спортсмены

В разделе **Спортсмены** тренер работает с карточками спортсменов.

Здесь можно:

- открыть список спортсменов;
- выбрать спортсмена;
- проверить профиль;
- обновить антропометрию, вес, пульс покоя, спортивную специализацию;
- указать сильные и слабые стороны;
- отметить ограничения, травмы и важные заметки;
- сохранить изменения профиля.

Рекомендуется заполнять профиль подробно: эти данные помогают точнее интерпретировать готовность, нагрузку и подготовку к стартам.

## 7. Планирование: общая логика

Раздел **Планирование** разделён по рабочему смыслу:

| Вкладка | Для чего нужна |
| --- | --- |
| План подготовки | Общая стратегия периода: цели, фазы, контроль, нагрузка, документ подготовки. |
| Мезоциклы | Блоки на несколько недель: фаза, прогрессия, целевая нагрузка, связь со стартом. |
| Недельный план | Основной рабочий экран тренера для сборки и назначения недели. |
| Библиотека шаблонов | Хранилище готовых дней, недель и полных планов. |
| Сезон и старты | Годовой календарь спортсмена, старты, план старта и результаты. |
| Календарь | Общий календарь соревнований и синхронизация с UWW. |

### Рекомендуемый порядок работы

1. Создайте или проверьте **Календарь** соревнований.
2. В **Сезон и старты** отметьте важные старты спортсмена.
3. В **План подготовки** задайте стратегию периода.
4. В **Мезоциклы** разбейте период на блоки.
5. В **Библиотека шаблонов** подготовьте или импортируйте шаблоны.
6. В **Недельный план** назначьте конкретную неделю спортсмену.
7. После выполнения смотрите **Review** и **Аналитику**.

## 8. План подготовки

Вкладка **План подготовки** не предназначена для ежедневного списка упражнений. Это стратегический документ.

Здесь указываются:

- название плана;
- период подготовки;
- описание цели;
- ключевые параметры;
- фазы подготовки;
- мониторинг состояния;
- общий документ подготовки.

Что сюда не нужно добавлять:

- отдельные упражнения на каждый день;
- назначение конкретной тренировки спортсмену;
- факт выполнения.

Эти действия выполняются в **Недельном плане** и **Библиотеке шаблонов**.

## 9. Мезоциклы

Мезоцикл - это блок подготовки на несколько недель.

Используйте вкладку **Мезоциклы**, чтобы:

- создать блок специальной, базовой, подводящей или другой подготовки;
- указать даты начала и конца;
- задать количество недель;
- выбрать фазу;
- связать блок со стартом;
- задать целевую нагрузку;
- отправить неделю из мезоцикла в недельный план.

Практический сценарий:

1. Выберите спортсмена в верхнем меню.
2. Откройте **Мезоциклы**.
3. Создайте блок на 3-6 недель.
4. Проверьте недели внутри блока.
5. Нажмите **В недельный план** на нужной неделе.
6. Перейдите во вкладку **Недельный план** и назначьте неделю спортсмену.

## 10. Недельный план

Это главный рабочий экран тренера.

Здесь тренер:

- видит рекомендованную неделю;
- проверяет предупреждения;
- применяет советы планировщика;
- смотрит ближайшие назначенные дни;
- задаёт дату начала недели;
- задаёт количество дней;
- выбирает плановую фазу;
- добавляет заметку тренера;
- назначает недельный план спортсмену.

Спортсмен в этой вкладке не выбирается отдельно. Используется спортсмен из верхнего меню.

### Как назначить недельный план

1. В верхнем меню выберите спортсмена.
2. Откройте **Планирование**.
3. Перейдите во вкладку **Недельный план**.
4. Укажите дату начала недели.
5. Укажите количество дней.
6. Проверьте фазу и заметку тренера.
7. Проверьте рекомендации и предупреждения.
8. Нажмите **Назначить недельный план**.

После назначения план появится у спортсмена в приложении и на сайте.

## 11. Библиотека шаблонов

**Библиотека шаблонов** - это хранилище заготовок. Это не список фактических назначений спортсмена.

Здесь можно:

- импортировать HTML-файл плана;
- сохранить план как один многодневный шаблон;
- открыть шаблон;
- редактировать день, сессию и упражнения;
- выбрать дни для назначения;
- назначить выбранный шаблон спортсмену;
- сохранить шаблон без назначения;
- удалить один или несколько шаблонов.

### Как импортировать план из файла

1. Откройте **Планирование**.
2. Перейдите в **Библиотека шаблонов**.
3. В блоке **Импорт плана работы** выберите HTML-файл.
4. Проверьте, что план загрузился как один полный шаблон.
5. Откройте нужный день в структуре плана.
6. При необходимости отредактируйте упражнения, подходы и контроль.
7. Нажмите **Сохранить в библиотеку** или **Назначить план**.

### Как удалить шаблоны

1. В списке шаблонов отметьте нужные чекбоксами.
2. Нажмите **Удалить выбранные**.
3. Подтвердите удаление.

Если шаблон уже назначен спортсмену, система может предупредить, что вместе с шаблоном будут удалены связанные назначенные тренировочные дни.

## 12. Сезон и старты

Вкладка **Сезон и старты** нужна для годового соревновательного плана спортсмена.

Здесь можно:

- создать сезон;
- выбрать год;
- отметить старты на календаре;
- выбрать режим отображения календаря;
- добавить план старта;
- сохранить результат старта;
- смотреть старты на годовой сетке.

Практический сценарий:

1. Выберите спортсмена в верхнем меню.
2. Откройте **Сезон и старты**.
3. Создайте сезон на нужный год.
4. Добавьте старты или перенесите их из календаря.
5. Для важного старта создайте план подготовки.
6. После соревнования заполните результат.

## 13. Календарь соревнований

Вкладка **Календарь** хранит общий список соревнований.

Здесь можно:

- добавить соревнование вручную;
- синхронизировать соревнования с UWW;
- фильтровать события по году, возрасту, стилю, типу и стране;
- выбрать соревнования;
- удалить выбранные соревнования;
- добавить соревнование в сезон спортсмена.

### Синхронизация с UWW

1. Откройте **Календарь**.
2. В блоке UWW выберите фильтры: год, возраст, стиль, тип, страна.
3. Нажмите **Обновить из UWW**.
4. Проверьте появившиеся соревнования.
5. Добавьте нужные старты в сезон спортсмена.

Календарь - это общий источник соревнований. Привязка к конкретному спортсмену выполняется через **Сезон и старты**.

## 14. Review

Раздел **Review** нужен для сравнения плана и факта.

Тренер смотрит:

- какие блоки выполнены;
- что спортсмен пропустил;
- где была изменена нагрузка;
- какие упражнения выполнены не по плану;
- какие есть отклонения по RPE, весу, подходам, объёму и длительности.

Используйте Review после тренировки или в конце недели.

## 15. Аналитика

Раздел **Аналитика** помогает принимать решения на основе данных.

Здесь можно смотреть:

- динамику готовности;
- тренд нагрузки;
- выполнение плана;
- риски накопления усталости;
- закономерности по спортсмену;
- тренерские рекомендации.

Аналитику лучше использовать после нескольких дней регулярного заполнения готовности и результатов. Чем больше данных внесено, тем полезнее выводы.

## 16. Работа спортсмена на сайте

Спортсмен видит свои рабочие зоны:

| Раздел | Что делает спортсмен |
| --- | --- |
| Готовность | Заполняет ежедневные данные перед тренировкой. |
| Тренировка | Смотрит назначенный день и отмечает выполнение. |
| История | Смотрит историю готовности и выполнения. |
| Соревнования | Видит ближайшие старты, фазу подготовки и дни до соревнования. |

## 17. Ежедневная готовность спортсмена

Спортсмен должен заполнять готовность каждый день до тренировки.

Обычно указываются:

- сон;
- качество сна;
- общее самочувствие;
- усталость;
- болезненность мышц;
- мотивация;
- пульс покоя;
- вес;
- боль;
- признаки болезни.

После отправки система показывает состояние дня. Цветовая логика:

| Статус | Значение |
| --- | --- |
| Green | Можно работать по плану. |
| Yellow | Нужен контроль или частичная адаптация. |
| Red | Высокий риск, нужна осторожность или изменение нагрузки. |

## 18. Тренировка спортсмена

В разделе **Тренировка** спортсмен видит назначенные дни и упражнения.

Что нужно сделать:

1. Открыть назначенный тренировочный день.
2. Посмотреть упражнения, подходы и контроль.
3. Выполнить работу.
4. Отметить выполненные упражнения.
5. Внести фактические значения, если они отличаются от плана.
6. Сохранить результат.

Если интернет пропал, результат сохраняется локально и будет отправлен позже при синхронизации.

## 19. Соревнования спортсмена

В разделе **Соревнования** спортсмен видит:

- ближайший старт;
- количество дней до старта;
- фазу подготовки;
- весовую цель;
- список соревнований;
- результат, если он уже внесён.

Если тренер меняет главный старт или план старта, у спортсмена обновляется количество дней до соревнования и соревновательный контекст.

## 20. Мобильное приложение

Мобильное приложение не открывает сайт внутри WebView. Интерфейс находится внутри приложения, а сервер используется только как API.

В приложении доступны:

- вход;
- dashboard;
- спортсмены для тренера;
- планы;
- календарь;
- готовность;
- результаты;
- отметка выполнения упражнений;
- офлайн-режим.

## 21. Офлайн-режим

Если нет интернета, приложение позволяет:

- открыть ранее загруженные данные;
- посмотреть назначенный план;
- заполнить готовность;
- внести результат тренировки;
- сохранить изменения локально.

Когда интернет появится, приложение отправит ожидающие действия на сервер.

Если синхронизация не прошла:

1. Проверьте интернет.
2. Откройте центр синхронизации.
3. Нажмите синхронизацию повторно.
4. Если ошибка повторяется, сообщите тренеру или администратору.

## 22. Типовой рабочий день

### Для спортсмена

1. Открыть приложение.
2. Заполнить готовность.
3. Посмотреть тренировку.
4. Выполнить упражнения.
5. Отметить, что сделано.
6. Внести фактические значения.
7. Сохранить результат.

### Для тренера

1. Открыть Dashboard.
2. Выбрать спортсмена.
3. Проверить готовность.
4. При необходимости посмотреть аналитику.
5. Проверить или скорректировать недельный план.
6. После тренировки открыть Review.
7. Принять решение по следующему дню.

## 23. Рекомендуемый недельный цикл тренера

1. В начале недели открыть **Недельный план**.
2. Проверить рекомендации и предупреждения.
3. Назначить недельный план.
4. Каждый день смотреть готовность спортсмена.
5. После тренировок смотреть выполнение.
6. В конце недели открыть **Review** и **Аналитику**.
7. На основе данных скорректировать следующую неделю.

## 24. Что делать при ошибках

| Ситуация | Что проверить |
| --- | --- |
| Не получается войти | Проверьте email, пароль и подключение к серверу. |
| Не виден спортсмен | Проверьте, выбран ли он в верхнем меню и привязан ли к тренеру. |
| Не виден назначенный план | Проверьте, что план действительно назначен выбранному спортсмену. |
| Не сохраняется результат | Проверьте интернет и очередь синхронизации. |
| Нет дней до старта | Проверьте, что у спортсмена есть активный старт в сезоне. |
| Не появились соревнования UWW | Проверьте фильтры календаря и повторите синхронизацию. |

## 25. Правила аккуратной работы

- Не создавайте дубли спортсменов.
- Не назначайте старые шаблоны без проверки содержимого.
- Перед удалением шаблонов проверяйте, назначены ли они спортсменам.
- Для каждого важного старта фиксируйте результат.
- Заполняйте готовность ежедневно, иначе аналитика будет неполной.
- После изменения плана проверяйте, как он выглядит у спортсмена.

## 26. Краткая схема платформы

```text
Календарь соревнований
  -> Сезон и старты спортсмена
  -> План подготовки
  -> Мезоциклы
  -> Недельный план
  -> Назначение спортсмену
  -> Готовность спортсмена
  -> Выполнение тренировки
  -> Review и аналитика
```

